% Phụ lục C: chi tiết cài đặt các cách giải thích đối chứng và các điểm sửa của bản đo v1. Không tính vào giới hạn trang.
% Nguồn: experiments/xai-metrics-v2/REPORT.md, experiments/xai-metrics/REPORT.md (mục đính chính), src/explain/xai_baselines.py.
\section{Chi tiết đánh giá lời giải thích}\label{app:xai}

\subsection{Cách chỉnh các cách giải thích GNN}
Các cách giải thích ở Mục~\ref{sec:xai-baselines} chạy trên cùng danh sách phần của ảnh hạng 1 (node vật thể và bộ ba có nhãn)
và chấm lại bằng chính reranker.
PGExplainer học một MLP sinh mask cho từng phần, trên 200 câu val nằm ngoài mẫu so sánh và ngoài 13 câu ở
Mục~\ref{sec:xai-cases}, với hệ số phạt của GNNExplainer. GNN-LRP lan truyền relevance theo quy tắc
$\varepsilon$~\cite{xai-bach2015lrp} qua các walk của hai lớp GAT, giữ chuẩn hoá và trọng số attention làm hằng số; điểm của
một phần là tổng relevance của các walk đi qua nó. GraphSVX ước lượng giá trị Shapley bằng hồi quy có trọng số (Kernel SHAP)
với cùng ngân sách lượt chấm như Shapley. SubgraphX tìm đồ thị con liên thông bằng MCTS (10 rollout, 10 mẫu cho mỗi ước lượng
Shapley cục bộ); thứ tự các phần theo đường cắt tỉa.

Các phép kiểm tra cài đặt: IG thoả completeness với sai số trung bình 0,7\% $\Delta$ (countonly) và 1,1\% (\texttt{main\_r3});
Shapley và GraphSVX thoả efficiency với sai số tối đa $1{,}3\cdot10^{-7}$; GNN-LRP dựng lại đầu ra GAT với sai số tối đa
$2{,}1\cdot10^{-6}$, và tổng relevance của các walk khớp relevance đầu ra, chỉ lệch do $\varepsilon$.

\subsection{Bản đo v1 sai ở đâu}\label{sec:xai-v1v2}
v1 được soát lại trước khi đưa vào báo cáo; nó giữ nguyên, kèm mục đính chính ở đầu \file{experiments/xai-metrics/REPORT.md}.
v2 sửa sáu điểm.
(1) Với \texttt{main\_r3}, v1 so term thô của hai kênh, dù fusion z-score cân hai kênh khác nhau; v2 thêm đạo hàm
ở~\eqref{eq:xai-term}. Sau khi đưa các term về cùng thang điểm, bộ ba đứng đầu E ở 70,5\% câu thay vì 1\%, và Fid$^{+}$@3
giảm từ \pmstd{65,7}{0,3}\% xuống \pmstd{34,9}{2,8}\%.
(2) v1 lấy lượt chấm thường làm gốc và tính câu E rỗng là ``không đổi''.
(3) Sparsity của v1 chia số dòng của bảng demo cho số phần của ảnh, hai đơn vị khác nhau.
(4) Tiêu chí ngữ nghĩa của v1 chấp nhận cả cặp khác loại; v2 giữ nó dưới tên lenient và thêm strict.
(5) v1 so khớp nhóm số lượng với \texttt{main\_r3} chấm trên đồ thị cấp 1 chưa gộp trùng, nơi hộp trùng làm tăng số đếm.
(6) v1 bỏ cạnh có nhãn là bỏ luôn bộ ba; đo lại dạng này cho 5,7\% câu mất hạng 1 (ngẫu nhiên: 5,1\%), so với 1,0\% khi chỉ
cắt truyền tin, nên hiệu ứng mà v1 thấy phần lớn đến từ việc mất bộ ba.
